\section{Scientific Interfaces and Operations}

Astronomical operations in Siderust combine typed quantities, epochs and
coordinate containers. A transformation receives the scientific context it
needs and returns a value with the resulting representation. The listings
below are included directly from executable sources, all checked against the
published v0.16.0 dependency graph during CI.

\subsection{Quantities and units}

\texttt{qtty} represents a quantity as \texttt{Quantity<U,S>}, where
\texttt{U} identifies the unit and \texttt{S} the scalar. Compatible
conversions are explicit, and multiplication or division constructs compound
units. This lets an application compute a velocity directly from a distance and an
elapsed time, with the compound unit carried by the result.

\lstinputlisting[float=tbp,caption={Computing a velocity from typed distance and duration quantities.},label={lst:qtty}]{../examples/src/bin/qtty.rs}

The local \texttt{Speed} alias names an existing velocity type; it introduces
neither a new nominal type nor an additional guarantee. We use such aliases
in applications to shorten notation while retaining the unit parameters.

Dimensional checking supports these operations automatically. Higher layers
add domain distinctions such as longitude and azimuth where their roles differ
scientifically, since a common angle unit alone does not distinguish them.

\subsection{Astronomical time}

\texttt{tempoch} represents an instant as a scale- and encoding-parameterized
\texttt{Time}. Subtracting two compatible instants produces a duration, and
shifting an instant by a duration produces another instant. JD, MJD, Unix, GPS,
and J2000-second markers are conversion targets rather than unlabelled numeric
epochs.

\lstinputlisting[float=tbp,caption={Typed time-scale and encoding conversion using \texttt{tempoch}.},label={lst:tempoch}]{../examples/src/bin/tempoch.rs}

UTC conversion uses the active UTC--TAI table. UT1 conversion requires an Earth-orientation or
$\Delta T$ context, whose source, interpolation policy and coverage interval
are part of the evaluated configuration \cite{iers2010}. The context exposes the available interval; callers of the evaluated EOP
lookup must additionally check its exact upper bound because of the endpoint
defect documented in Section~8.

\subsection{Coordinates, frames, and centers}

\texttt{affn} defines \texttt{Position<C,F,U>}, where \texttt{C} is the
reference center, \texttt{F} the frame, and \texttt{U} the length unit.
A position transformation changes its frame or center marker explicitly.
Directions carry orientation independently of an origin, making them suitable
for distance-free pointing operations. Incompatible markers are checked when
values are combined.

Reference-center types may carry runtime parameters. In the evaluated snapshot,
\texttt{Topocentric} uses a \texttt{Geodetic<ECEF>} site as its parameter type.
\texttt{Bodycentric} instead uses a \texttt{BodycentricParams} value containing
an elliptic orbit, its reference center, and a gravitational parameter.
Coordinates with the same center type can therefore still carry different
parameters. Cartesian position subtraction is available as an operator only for
centers whose parameter type is \texttt{()}. For parameterized centers,
\texttt{checked\_sub} returns an error when the stored values differ.
The accompanying tests exercise both matching and mismatched sites.
Transform methods preserve or explicitly replace the
stored parameters. Site values identify the geometric observer; external station metadata and
site epochs are managed by the application.

\subsection{Generic surface observers}
\label{sec:surface}

We generalize surface observers through
\texttt{SurfaceSite<B>}. A body model associates \texttt{B} with a center,
body-fixed frame and default reference surface. The same transformation
infrastructure then serves lunar, terrestrial and other supported surfaces.
The body parameter is static; longitude, latitude, height, surface parameters
and their admissibility are checked values. Planetographic latitude identifies
the ellipsoid normal, whereas planetocentric latitude selects a radial
intersection with the surface. Height and local vertical follow the surface
normal in both cases; on a sphere, these latitude definitions coincide.

A finite position is translated to the body's center and the observer's
location is subtracted before local angular extraction, retaining range and
parallax. A distance-free direction undergoes rotation only. The local east--north--up
basis is right-handed; the public \texttt{Horizontal} Cartesian ordering is
north, east, up, consistent with azimuth measured from north through east.
The component order is explicit at the public interface.

The high-level direction observation retains both \texttt{B} and the site
values. Checked inversion rejects a different site, while a different body
is rejected by the type system. Low-level \texttt{Direction<Horizontal>}
values do not themselves retain the observer. The epoch and orientation context are supplied explicitly to each observation.
This design connects static body identity with value-dependent observer geometry.

\subsection{Astronomy workflows}

The domain crate supplies transforms, ephemerides, bodies, observatories,
altitude calculations, observing-event searches, and feature-gated dynamics
workflows. Analytical planetary and lunar paths use archive-backed VSOP87 and
ELP2000-82B data \cite{bretagnon1988vsop87,chapront1983elp2000}; runtime paths
load externally supplied JPL/NAIF BSP kernels \cite{naifspk,park2021de440}.

\lstinputlisting[float=tbp,caption={Runtime SPK-backed ephemeris access. The application supplies the BSP file at execution time.},label={lst:runtime-eph}]{../examples/src/bin/runtime_ephemeris.rs}

The common ephemeris interface lets an application select an analytical model
or a supplied SPK kernel through common frame and center types. Their physical
frame realizations are specified in Section~7.1.
The evaluation records that choice together with the body definition and
correction mode, connecting each result to its scientific model.


\subsection{Three levels of scientific contract}

Static contracts constrain compatible unit, frame, center, scale and observer-body
parameters. Dynamic contracts validate values sharing those types: two sites
can have different coordinates, a vector can be degenerate, and a query can exceed
a dataset's coverage. Scientific contracts concern the intended observable:
ephemeris choice, frame realization, geometric or apparent corrections, body
orientation, surface definition and data accuracy. These require model selection
and numerical validation rather than type checking.

Table~\ref{tab:guarantees} summarizes the executed static and dynamic controls.
The semantic matrix below summarizes the observed enforcement mechanisms and
outcomes for the evaluated public APIs.
The evaluation included five compiler-failure controls; runtime cases covered
nonfinite values, invalid latitude, a surface-normal singularity, zero direction
and mismatched observers. Dataset coverage was examined separately, revealing
the EOP and exact SPK endpoint behaviors described in the results.

\begin{table*}[!tbp]
\centering
\caption{Static and runtime semantic controls observed with published v0.16.0.
Scientific model adequacy is assessed by numerical references, not by these contracts.}
\label{tab:guarantees}
\footnotesize
\begin{tabularx}{\textwidth}{@{}YYYY@{}}
\toprule
Constraint & API & Mechanism & Outcome \\
\midrule
observer\_\allowbreak{}body & horizontal\_\allowbreak{}observation\_\allowbreak{}to\_\allowbreak{}direction$<$B$>$ & static type/trait compatibility & rejected: E0308 \\
center & Position$<$C,F,U$>$ subtraction & static type/trait compatibility & rejected: E0277 \\
dimension & Quantity$<$U,S$>$ addition & static type/trait compatibility & rejected: E0308 \\
frame & Position$<$C,F,U$>$ subtraction & static type/trait compatibility & rejected: E0277 \\
scale & gmst\_\allowbreak{}iau2006(UT1,TT) & static type/trait compatibility & rejected: E0308 \\
observer values & Position::checked\_\allowbreak{}sub & runtime Result & valid accepted; invalid rejected \\
finite inputs / surface domain & SurfaceSite::new & runtime Result & valid accepted; invalid rejected \\
zero/nonfinite direction & SurfaceHorizontalExt::observe\_\allowbreak{}direction & runtime Result & valid accepted; invalid rejected \\
surface observer values & horizontal\_\allowbreak{}observation\_\allowbreak{}to\_\allowbreak{}direction & runtime Result & valid accepted; invalid rejected \\
\bottomrule
\end{tabularx}

\end{table*}
